\section{Introduction}\label{sec:intro}

\subsection{The problem}

Many social processes rest on a bank of test items: civic-knowledge questionnaires,
professional certifications, licensing examinations, the claims checked by a fact-checking
service. Whoever writes and selects the items shapes what the test measures, and in many
settings no institution is credibly impartial. A government that selects the questions of a
civic test can select questions that flatter it. A professional body that writes its own
licensing exam has a stake in who passes it.

The obvious decentralized alternative is to let participants propose items and vote on them.
Voting moves control from the institution to whichever faction holds the majority. Reputation
systems in which approval earns weight make matters worse: a cohesive group gains weight
faster than a dispersed one and then uses it. The hard part is not collecting items. It is
deciding which items are fair \emph{without trusting a single authority and without turning
the decision into a popularity contest}.

\subsection{The mechanism in brief}

Isegoria\footnote{Greek \emph{isēgoria}: the equal right of every citizen to speak in the
assembly.} is a design for this problem~\cite{isegoria-repo}. It rests on three pillars, in
order of importance.
\begin{enumerate}
  \item \textbf{Uniqueness without identification.} Participants are anonymous, but each
    anonymous participant is backed by exactly one real person. Uniqueness is certified by an
    external identity system and then cryptographically severed from everything the person
    does. Without it no aggregation rule survives a Sybil attack~\cite{douceur2002sybil}.
  \item \textbf{A quality signal that is not an opinion.} Items are trialled on samples of
    respondents, and the final verdict comes from psychometrics on the response data: item
    response theory (IRT) for discrimination and difficulty, and differential item
    functioning (DIF) for bias~\cite{lord1980,holland1993dif}. Because no attribute of any
    person is ever collected, DIF is detected on \emph{latent} classes.
  \item \textbf{Bridging instead of majority.} Peer review is aggregated by a matrix
    factorization that credits approval only to the extent that it is not explained by a
    latent axis of division, in the style of Community Notes~\cite{wojcik2022birdwatch}. An
    item approved only by the largest camp is filtered out.
\end{enumerate}
Peer review (Level~A) is only an upstream filter deciding which items are worth trialling. The
verdict is empirical (Level~B). Reviewer reputation (Level~C) weights the review and is scored
against outcomes, with a proper scoring rule so that honest forecasting is the best strategy. A
sublinear discount limits coordinated blocs. Because bridging cannot tell a true but divisive
fact from partisan propaganda, an author may appeal a polarization rejection directly to the
evidence filter, staking reputation. No money enters the system at any point.

\subsection{Contributions}

This paper states the mechanism formally and analyses it. The analysis is the main
contribution. Several of its results expose properties that the design documents do not
anticipate. We report them as findings with proposed corrections, not as settled
features of the design.

\begin{enumerate}
  \item \textbf{Formal statement} (\cref{sec:setting}). Requirements, adversaries and a set
    of assumptions that abstract the identity and network layers. The analysis therefore
    holds for any realization of those layers that meets the assumptions.
  \item \textbf{Level A} (\cref{sec:levela}). We characterize the invariances of the
    bridging objective (\cref{lem:gauge}) and derive three consequences. Bridge scores sum to
    zero within a fit, so the gate is relative to the batch (\cref{prop:zerosum}). The origin
    of the latent axis, which decides whether the score is camp-balanced or majoritarian, is
    fixed only by the regularization; in a two-camp model we obtain the interpolation weight
    in closed form (\cref{prop:leak}) and confirm it on the full model (\cref{tab:leak}).
    Finally, no rule that reads only ratings can separate true-but-divisive items from
    partisan ones (\cref{prop:equivalence}).
  \item \textbf{Level B} (\cref{sec:levelb}). We show why latent-class DIF is identified by
    the conditional covariance between biased items, and why a single biased item is nearly
    invisible (\cref{prop:single,prop:pair}). We also show that measurement error in the
    ability proxy produces spurious latent classes (\cref{prop:proxy}), and we measure its
    effect on the production detector.
  \item \textbf{Level C} (\cref{sec:levelc}). The ratio-form Brier skill score is not proper,
    a known property of skill scores that here reverses the intended incentive (\cref{prop:bss}). A leave-one-out difference score is proper and keeps the property that
    matters for the design, namely zero reward for following the crowd (\cref{prop:loo}).
    As currently parameterized, the weight cap never binds (\cref{prop:cap}).
  \item \textbf{Coordination} (\cref{sec:collusion}). The sublinear discount rewards evading
    its detector (\cref{prop:split}). At design scale reviewers share almost no items within
    an epoch (\cref{prop:overlap}), so the per-epoch defence is random assignment.
  \item \textbf{Evidence and plan} (\cref{sec:implementation,sec:limitations,sec:plan}).
    The evidence level of every claim, the structural limits, and an evaluation plan with
    decision rules fixed in advance.
\end{enumerate}

\subsection{Status and reading guide}

This is a design-and-analysis paper about a research-stage system. It does not claim that the
mechanism works in the field. No component has been externally reviewed, every operational
threshold is provisional, and the identity and network layers are partly scaffolds. To keep
claims honest, results carry one of four labels, following the verification methodology of
the repository~\cite{isegoria-repo}:
\proved{} (a proof is given),
\empirical{} (a simulation result, valid in the stated regime only),
\designchoice{} (a choice that is motivated but not derived), and
\openq{} (a question we cannot yet answer).
Every number in the paper can be regenerated with the scripts listed in
\cref{app:reproducibility}. Where we report the behaviour of the production code, we run the
Rust engine itself rather than a re-implementation.
